\begin{figure}[tb]
\centering
\begin{tikzpicture}
\begin{axis}[diagnostic axis,title={Reflect the same curve},ylabel={$f(y)$},
 xmin=-1,xmax=1,ymin=.5,ymax=1.6]
 \addplot[reference quantity,strong line,domain=-1:1]{1+.2*x^2+.3*x};
 \addlegendentry{$f(y)$}
 \addplot[estimated quantity,alternative pattern,standard line,domain=-1:1]{1+.2*x^2-.3*x};
 \addlegendentry{$f(-y)$}
\end{axis}
\begin{axis}[diagnostic axis,at={(.5\textwidth,0)},anchor=south west,
 title={Average and half-difference},ylabel={Component},xmin=-1,xmax=1,ymin=-.4,ymax=1.4]
 \addplot[derived quantity,strong line,domain=-1:1]{1+.2*x^2};
 \addlegendentry{$\evenq[f]$}
 \addplot[torque quantity,standard line,domain=-1:1]{.3*x};
 \addlegendentry{$\oddq[f]$}
\end{axis}
\end{tikzpicture}
\caption{Reflection separates an illustrative curve
$f(y)=1+0.2y^2+0.3y$ into its allowed even part and its odd tilt.
The same projections applied to d and q flux select opposite diagnostic
channels. The schematic curve is not a second machine model.}
\label{fig:decomposition}
\end{figure}
